\section{Experiments}
\label{sec:experiments}

\subsection{Experimental settings}
\label{sec:setup}
\subsubsection{Benchmark instances}
\label{sec:setup:instances}

Instances come from the MKT benchmark~\citep{homayouni2021}, which takes the
processing data of the Brandimarte MK instances~\citep{brandimarte1993}
unchanged and adds two things: a machine-to-machine travel-time matrix and a
vehicle count. The matrices are asymmetric, and their zeroth row and column
address the load--unload station. Adding transport to the MK data does not
alter any processing time, so results can be traced back to the original
instances.

This paper reports results on MK01 only: 10 jobs, 6 machines, 55 operations.
We state this plainly because it bounds every claim below. The remaining
instances of the benchmark are not covered here, and no statement in this paper
should be read as extending to them.

One property of MK01 matters for the interpretation of the results. Of the
instances in the benchmark, MK01 offers the fewest opportunities for
multi-drop travel: the instance can batch more than one task on 1.4\% of its
pick-ups, against 2.5\% for MK07 and 12.3\% for MK10. A mechanism that only
acts when a batch can be formed will therefore look inert on MK01 for reasons
that belong to the instance rather than to the mechanism. We report the
batching results with that caveat attached rather than as a verdict on the
mechanism.

\subsubsection{Configuration}
\label{sec:setup:config}

Two report tiers are used throughout. Tier A (\texttt{A-MKT}) switches off all
ten constraints and all mechanisms, leaving a flexible job shop with transport.
Tier B (\texttt{B-Full}) enables all of them. The two tiers differ only in the
mechanism switches; the instance, the fleet size, and the transport caliber are
identical.

All policies are trained for 300 steps with a group size of $G = 8$, learning
rate $3\times10^{-4}$, and Adam, on a single training seed ($\text{seed} = 0$).
The policy is evaluated every 30 steps on five evaluation seeds with argmax
action selection. The training seed is the only source of variation in the
learned parameters; the five evaluation seeds perturb the simulation, not the
training.

We report the training-seed variation explicitly, because it is large. Under
the \texttt{scalar} advantage with the budget mechanism off, five training
seeds on MK01 give makespans of 84.04, 81.24, 87.19, 79.27, and 71.56 --- a
mean of 80.66, a standard deviation of 5.90, and a range of 15.63. Every
ablation number in Section~\ref{sec:results} other than the advantage ablation
comes from a single training seed, and the differences between ablation arms
are of the same order as that standard deviation. Those numbers are reported as
measurements, not as effect estimates, and are labelled as single-training-seed
throughout. The advantage ablation (Section~\ref{sec:results:advantage}) is the
exception: all three of its arms were trained with five seeds, precisely because
its conclusions are not readable from one.

\subsubsection{Baselines}
\label{sec:setup:baselines}

Three baselines are used.

\begin{itemize}
\item \textbf{Rule.} Each operation goes to its fastest eligible machine, with
      a first-in-first-out tie-break, and vehicles are dispatched round-robin.
      This is also the schedule used to derive the reward weights of
      Equation~\eqref{eq:weights}, so it defines the reference point that the
      reward is normalized against.
\item \textbf{Single-objective GRPO.} The same policy and training procedure,
      with the preference vector set to a one-hot vector over one of the three
      objectives. Three arms are trained, one per objective. This is the
      standard scalarization control: it answers whether training on the
      weighted scalar is better or worse than training once per objective and
      choosing afterwards.
\item \textbf{NSGA-II.} Population 100, 50 generations, 5000 evaluations per
      tier, run on the same instance and the same fleet and constraint
      configuration as the learned policies.
\end{itemize}

A fourth baseline layer --- a published deep reinforcement learning method ---
is not included. The comparison is therefore restricted to a rule, a
single-objective version of our own method, and a metaheuristic.
Section~\ref{sec:discussion} records what this omission costs.

\paragraph{Out-of-sample re-evaluation of the front}
NSGA-II searches on a single random chain. Points found on that chain can
overfit it. Every point of the NSGA-II front is therefore re-evaluated on the
same five evaluation seeds used for the policies, and the re-evaluated front is
what we compare against. The difference is not cosmetic: on tier B the search
front reads $78.09$ in-sample and $79.56$ after re-evaluation, and individual
front points move further ($78.18 \rightarrow 88.62$).

\paragraph{Matching the baselines to the tier}
The rule baseline is run under the same constraint set as the policy it is
compared with, using the same five evaluation seeds, so that comparisons
between the rule row and the policy row are paired. The rule makespan under
tier B (105.35) and under tier A (73.01) are therefore different quantities,
and neither equals the reference makespan used for the reward weights.

\subsubsection{Reproducibility}
\label{sec:setup:repro}

Every mechanism switch defaults to off, and a run with all switches off
reproduces the pre-mechanism trajectory bit for bit. This property is what
makes the ablations in Section~\ref{sec:results} measurements of the mechanism
rather than of an incidental change elsewhere, so it is verified rather than
asserted: a digest of the sampled chain and of the parameter vector is pinned
for each switch position, and the same discipline is applied to the wiring
between training and evaluation, so that a switch cannot be active in one and
inactive in the other.

Because the switches alter the shop dynamics, the reference makespan differs
between ablation arms. Disabling age-dependent failure changes the reference
makespan from 109.05 to 117.77; disabling multi-drop travel changes it to
115.13. The absolute makespan of an arm that switches off a dynamics constraint
is therefore not comparable with the absolute makespan of an arm that leaves it
on: the two arms solve different problems. Where such a comparison would be
misleading we report the paired difference against the all-on configuration
instead.

\subsubsection{Why published figures are not reproduced in our tables}
\label{sec:related:incomparable}

We deliberately do not place our numbers in a table beside published numbers
from this literature. Four differences independently prevent the comparison,
and each one is large enough to change the conclusion on its own.

\emph{Travel-time caliber.} Some published results use a machine-to-machine
travel-time matrix and others use geometric travel times. We use both, and the
caliber is a property of the instance rather than a run-time switch, so a
number cannot be moved between calibers by changing a configuration. The same
chromosome evaluated under a geometric caliber and under a matrix caliber
differs by roughly a factor of four in our own measurements.

\emph{Fleet size.} Our main configuration uses three vehicles. Published
results on the same instance family use six or use the number of machines. The
vehicle count is not published with the benchmark: the source
paper~\citep{homayouni2021} describes it as a uniform draw over a range of the
machine count, so any specific value used by a later paper is a reconstruction
rather than a reproduction.

\emph{Simulator.} Each study runs its own discrete-event simulation. Two
simulators that agree on the problem statement can disagree on makespan through
tie-breaking, event ordering, and rounding.

\emph{Transport boundary.} Our makespan includes the travel of a job from the
load--unload station to its first machine and its return travel after its last
operation. Published makespans in this line end at the completion of the last
operation. The difference is measurable and not small.

Because these four differences are independent, no single adjustment corrects
for them all, and a table that placed the numbers side by side would invite a
comparison that the numbers cannot support. We therefore compare only against
baselines we ran ourselves on our own configuration, and we state the boundary
explicitly rather than leaving it to a footnote. The consequence for evaluation
is that this paper cannot claim to have beaten a published method; it can claim
to have measured its own method against a rule, a single-objective version of
itself, and a metaheuristic, all under one caliber.
Section~\ref{sec:discussion} records what that costs.

\label{sec:results}

\subsection{Main comparison}
\label{sec:results:main}

Table~\ref{tab:main} reports the three baselines against the learned policy on
both tiers. The $\pm$ values are the spread over the five evaluation seeds.

\begin{table*}[htbp]
\centering
\caption{Main comparison on MK01. Tier A disables all ten constraints and all
mechanisms; tier B enables them. The
rule row is run under the same constraint set as the policy it is paired with,
so the two rule rows are not the same quantity. NSGA-II rows are re-evaluated
out-of-sample on the five evaluation seeds.}
\label{tab:main}
\setlength{\tabcolsep}{5pt}
\begin{threeparttable}
\begin{tabular}{@{}llrrrr@{}}
\toprule
Tier & Method & Makespan & Energy (\kWh) & \TWT & Source \\
\midrule
A & Rule (same constraint set) & 73.01 $\pm$ 0.00 & 5.675 & 0.000 & --- \\
A & Single-objective GRPO, makespan & 57.60 & 4.714 & 0.000 & --- \\
A & Single-objective GRPO, energy & 52.91 & 4.341 & 0.000 & --- \\
A & \textbf{Ours} (weighted scalar) & 49.27 & 4.097 & 0.000 & --- \\
A & NSGA-II (out-of-sample) & \textbf{49.14} & \textbf{4.093} & 0.000 & --- \\
\midrule
B & Rule (same constraint set) & 105.35 $\pm$ 4.56 & 7.936 $\pm$ 0.33 & 60.14 $\pm$ 19.00 & --- \\
B & Single-objective GRPO, makespan & \textbf{74.09} $\pm$ 4.67 & 5.790 & 5.977 & --- \\
B & Single-objective GRPO, energy & 96.93 $\pm$ 9.98 & 7.434 & 97.458 & --- \\
B & Single-objective GRPO, \TWT & 77.30 $\pm$ 6.35 & 6.048 & 7.392 & --- \\
B & \textbf{Ours} (weighted scalar) & 78.90 $\pm$ 3.71 & \textbf{6.160} & \textbf{5.042} & --- \\
B & NSGA-II (out-of-sample) & 79.56 $\pm$ 2.64 & 6.239 & 42.48 & --- \\
\bottomrule
\end{tabular}
\begin{tablenotes}\footnotesize
\item Tier A switches off the tardiness constraint, so \TWT\ is identically
      zero for every row of tier A by definition, not as a measurement. Tier A
      also has no stochastic source left once the ten constraints are off, so
      its five evaluation seeds return identical values and no statistical test
      is possible on that tier. The tier A NSGA-II row is deterministic and
      rests on 5000 evaluations.
\end{tablenotes}
\end{threeparttable}
\end{table*}



On tier B the learned policy dominates the entire NSGA-II front: the front's
best re-evaluated values are 79.56 for makespan, 6.239 for energy, and 20.73
for tardiness, and the policy attains 78.90, 6.160, and 5.042. Every one of the
12 re-evaluated front points is dominated. The rule baseline is beaten by
26.5 in makespan and by 55.1 in tardiness.

On tier A the ordering reverses, narrowly. NSGA-II reaches 49.14 against our
49.27 and 4.093 against our 4.097. A difference of 0.3\% in makespan on a
deterministic problem is not a tie in the statistical sense --- there is no
noise to attribute it to --- and we report it as a loss rather than rounding it
away. The reading is that on the easier problem, where the ten constraints are
absent and 5000 deterministic evaluations are available, the metaheuristic is
at least as good.

The single-objective arms do not displace the weighted scalar. On tier B the
makespan-only arm achieves the best makespan of any learned row (74.09) but
pays for it in tardiness (5.977 against 5.042), and the energy-only arm is
worse on every objective. The spread across the last few evaluations of a single
run is 10--20\% of the value, which is larger than most of the gaps here.

\subsection{Constraint and mechanism ablation}
\label{sec:results:ablation}

Table~\ref{tab:ablation} reports twelve ablation arms against the all-on
configuration. Each arm switches off one constraint group or one mechanism.

\begin{table*}[htbp]
\centering
\caption{Ablation against the all-on configuration (\texttt{full}, makespan
78.90). $\Delta$ is arm minus \texttt{full}: positive means the arm is worse,
so a positive $\Delta$ is the mechanism helping. $p$ values are paired
$t$-tests over five evaluation seeds against \texttt{full}.}
\label{tab:ablation}
\setlength{\tabcolsep}{5pt}
\begin{threeparttable}
\begin{tabular}{@{}lrrrrrr@{}}
\toprule
Arm & $\Delta$makespan & $p$ & $\Delta$energy & $p$ & $\Delta$\TWT & $p$ \\
\midrule
\texttt{c-no-logistics}  & $+1.69$ & 0.453 & $+0.116$ & 0.481 & $+19.07$ & \textbf{0.012} \\
\texttt{c-no-production} & $-8.17$ & \textbf{0.012} & $-0.619$ & \textbf{0.009} & $-2.65$ & 0.535 \\
\texttt{c-no-info}       & $+5.41$ & \textbf{0.028} & $+0.406$ & \textbf{0.024} & $+28.06$ & \textbf{0.020} \\
\texttt{c-none}          & $-29.63$ & \textbf{0.0001} & $-2.063$ & \textbf{0.0001} & $-5.04$ & 0.267 \\
\texttt{m-no-route}      & $-1.64$ & 0.581 & $-0.075$ & 0.729 & $+6.58$ & 0.448 \\
\texttt{m-no-r2}         & $+7.79$ & \textbf{0.048} & $+0.531$ & 0.065 & $+30.19$ & 0.068 \\
\texttt{m-no-geom}       & $+8.28$ & \textbf{0.001} & $+0.664$ & \textbf{0.001} & $+26.23$ & 0.077 \\
\texttt{m-no-pm}         & $+1.30$ & 0.146 & $+0.103$ & 0.130 & $+0.49$ & 0.931 \\
\texttt{m-no-charge}     & $-1.15$ & 0.591 & $-0.062$ & 0.677 & $+4.08$ & 0.616 \\
\texttt{m-no-batch}      & $+5.66$ & \textbf{0.008} & $+0.376$ & \textbf{0.013} & $+34.74$ & \textbf{0.018} \\
\texttt{m-no-t3}         & $+5.14$ & 0.187 & $+0.425$ & 0.142 & $+16.39$ & 0.056 \\
\texttt{adv-reinforce}   & $-7.78$ & 0.053 & $-0.587$ & \textbf{0.049} & $-3.83$ & 0.434 \\
\bottomrule
\end{tabular}
\begin{tablenotes}\footnotesize
\item Arms beginning with \texttt{c-} switch off constraint groups; arms
      beginning with \texttt{m-} switch off a single mechanism. Six arms
      necessarily disable more than the named item, because a mechanism cannot
      be enabled when the constraint it acts on is off; the released
      configuration files record the exact switch set of every arm. The
      \texttt{adv-reinforce} row here is the single-seed run; the five-seed
      version is in Table~\ref{tab:advantage}.
\end{tablenotes}
\end{threeparttable}
\end{table*}

\begin{figure*}[htbp]
\centering
\includegraphics[width=\textwidth]{fig_ablation.pdf}
\caption{Ablation against the all-on configuration, the same numbers as
Table~\ref{tab:ablation}. The shaded band is one training-seed standard
deviation. The tardiness differences mostly fall outside it; the makespan
differences mostly do not.}
\label{fig:ablation}
\end{figure*}

Table~\ref{tab:ablation}, plotted in Figure~\ref{fig:ablation}, reads as two
groups.

Among the mechanism arms, the separation is between the mechanisms the
instance can exercise and the mechanisms it cannot. Three of them have nothing
to act on here: the charging constraint never activates under these parameters,
the maintenance interval is never reached on this instance, and congestion does
not bind at the reference load. Their makespan differences straddle zero
($-1.6$ to $+1.3$) and their tardiness differences are small ($+0.5$ to
$+6.6$), which is what a correctly implemented mechanism with no effect to make
should look like. The four mechanisms that do have something to act on --- the
aisle-segment tokens, the distance bias, the batch head, and the budget
mechanism --- have positive differences in every case: $+5.1$ to $+8.3$ in
makespan and $+16$ to $+35$ in tardiness. The grouping is not chosen after the
fact. Whether a mechanism can act is established independently of this table,
by measuring the constraint rather than the policy
(Tables~\ref{tab:active} and~\ref{tab:batching}).

What the table does not support is an ordering among the mechanisms. With one
training seed per arm, and a measured training-seed standard deviation of
5.90, the differences are of the same order as the seed-to-seed variation, and
a re-run with other seeds would shuffle the order. A single seed cannot
separate a mechanism that helps by five makespan units from a mechanism that
does nothing. The $p$ values in the table test consistency across five
\emph{evaluation} seeds, which is a different source of variation and does not
cover the training seed. We therefore report the grouping and not the ordering,
and Section~\ref{sec:conclusion} states how many seeds an ordering would need.
The three negative arms among the constraint groups --- production, all-off,
and the REINFORCE advantage --- fall inside the range that changing the
training seed alone can produce.

One further caveat applies to the three constraint-group arms. Switching off a
dynamics constraint changes the problem: the reference makespan moves from
109.05 to 117.77 when age-dependent failure is off and to 115.13 when multi-drop
travel is off. The absolute makespans of those arms are therefore not
comparable to the all-on configuration, and only the paired differences above
are meaningful.

The levels behind those differences are worth stating, because the differences
are differences of small numbers. Across the thirteen arms the total energy
account spans $0.0772$ to $0.0832$ kWh per minute of makespan, a range of
7.6\% of its mean, while the net account spans four times as much; machine
utilization ranges from $0.399$ to $0.704$, and the highest value belongs to
the all-off arm, which finishes ten jobs quickly because it faces the easiest
problem.



\subsection{What the benchmark actually exercises}
\label{sec:results:active}

A constraint that is modelled but never triggered contributes nothing to the
problem, and whether a constraint triggers is an empirical question. It can be
answered without training: run the reference schedule and count the events.
This measurement is a property of the constraint under the reference schedule,
not a property of any learned policy, so we report it for all ten instances of
the benchmark even though every policy result in this paper comes from MK01.

\begin{table}[htbp]
\centering
\caption{Constraint activation under the reference schedule, tier B switches,
one seed. Counts are per episode. The maintenance clock column gives the
largest accumulated spindle time reached by any machine, against a maintenance
interval of 120 minutes.}
\label{tab:active}
\setlength{\tabcolsep}{2.5pt}
\begin{threeparttable}
\begin{tabular}{@{}lrrrrr@{}}
\toprule
Instance & $C_{\max}$ & Failures & Forced PM & Max clock & Drained \\
\midrule
mk01 & 109.05 & 0 & \textbf{0} & \textbf{82} & 0 \\
mk02 & 74.72 & 0 & \textbf{0} & \textbf{45} & 0 \\
mk03 & 406.35 & 1 & 3 & 125 & 0 \\
mk04 & 270.75 & 0 & 1 & 122 & 0 \\
mk05 & 351.91 & 1 & 4 & 126 & 0 \\
mk06 & 233.77 & 1 & \textbf{0} & \textbf{105} & 0 \\
mk07 & 309.72 & 0 & 3 & 128 & 0 \\
mk08 & 777.85 & 5 & 16 & 133 & 0 \\
mk09 & 616.91 & 1 & 14 & 129 & 0 \\
mk10 & 501.02 & 1 & 9 & 128 & 0 \\
\bottomrule
\end{tabular}
\begin{tablenotes}\footnotesize
\item No vehicle ever drains on any instance, so the charging constraint is
      inert throughout the benchmark under the standard parameters. The
      maintenance constraint activates on seven of the ten instances; on the
      three marked in bold the accumulated spindle time stays below the
      interval, so no maintenance is ever forced. Vehicle failures occur 15 to
      114 times per episode depending on the instance and are not tabulated
      here.
\end{tablenotes}
\end{threeparttable}
\end{table}

Two consequences run through the rest of the paper. First, the charging
mechanism has nothing to act on: the fleet consumes 0.51--4.62 \kWh{} per
episode against a 2--4 \kWh{} battery, a factor of three to ten of headroom, so
the state never arises. Second, the maintenance mechanism has nothing to act on
in MK01 specifically, which is the instance all policy results come from. Both
are reported as facts about the parameterization rather than as properties of
the mechanisms.

\subsubsection{Where congestion begins to bind}
\label{sec:results:load}

The congestion mechanism shows a zero in Table~\ref{tab:ablation} on MK01, and
it is worth establishing whether that belongs to the instance or to the
parameters. Table~\ref{tab:load} varies the travel load through the vehicle
speed and the supply through the fleet size, and records the time vehicles
spend waiting for an aisle segment as a share of the makespan per vehicle.

\begin{table}[htbp]
\centering
\caption{Aisle contention as a function of travel load and fleet size, under
the reference schedule with all ten constraints on. Entries are the total time
vehicles wait for an aisle segment, divided by the fleet size and by the
makespan, in per cent.}
\label{tab:load}
\setlength{\tabcolsep}{5pt}
\begin{threeparttable}
\begin{tabular}{@{}lrcc@{}}
\toprule
Instance & Vehicle speed (m/s) & 3 vehicles & 6 vehicles \\
\midrule
mk01 & 1.00 & 0.31 & 0.20 \\
mk01 & 0.50 & 0.51 & 0.65 \\
mk01 & 0.25 & 1.23 & 1.33 \\
mk01 & 0.10 & \textbf{18.13} & \textbf{10.41} \\
\midrule
mk10 & 1.00 & 0.13 & 0.12 \\
mk10 & 0.50 & 0.27 & 0.35 \\
mk10 & 0.25 & 0.97 & 0.90 \\
mk10 & 0.10 & \textbf{4.94} & \textbf{5.09} \\
\bottomrule
\end{tabular}
\begin{tablenotes}\footnotesize
\item The nominal speed is 0.5~m/s. With a single vehicle the entry is zero at
      every speed, since one vehicle cannot contend with itself.
\end{tablenotes}
\end{threeparttable}
\end{table}

Three things are readable. First, the single-vehicle column is identically
zero at every speed, which is a property of the metric rather than a finding:
it confirms that what is being measured is interaction between vehicles and not
travel time as such. Second, at the nominal speed a vehicle waits for an aisle
segment well under one per cent of the time, and adding vehicles does not
change that. Third, the transition to a contended regime is abrupt rather than
gradual: on MK01 the wait goes from 1.2\% at a quarter of the nominal speed to
18\% at a tenth of it, and on MK10 from 1.0\% to 4.9\%. The constraint is
therefore not inert in general. It is inert at the benchmark's operating point,
where vehicle utilisation lies between 4\% and 23\%, and that is the sense in
which its zero in Table~\ref{tab:ablation} should be read.

\subsubsection{Batching opportunities}
\label{sec:results:batching}

The multi-drop travel model and its batching head were added because the load
capacity of a vehicle was a parameter the model carried but never used.
Whether the mechanism has anything to do depends on how often a pick-up point
holds more than one waiting task. Table~\ref{tab:batching} measures that
across the three instance sizes, again under the reference schedule.

\begin{table}[htbp]
\centering
\caption{Batching opportunities under the reference schedule. ``Pick-ups'' is
the number of times a vehicle collects at a pick-up point; the middle column
counts how often more than one task was available there.}
\label{tab:batching}
\setlength{\tabcolsep}{4pt}
\begin{threeparttable}
\begin{tabular}{@{}lrrrr@{}}
\toprule
Instance & Seed & Pick-ups & Multi-task & Mean per trip \\
\midrule
mk01 & 0 & 70 & 1 (1.4\%) & 1.014 \\
mk01 & 1 & 67 & 1 (1.5\%) & 1.015 \\
mk07 & 0 & 119 & 3 (2.5\%) & 1.042 \\
mk07 & 1 & 120 & 2 (1.7\%) & 1.033 \\
mk10 & 0 & 235 & 29 (12.3\%) & 1.140 \\
mk10 & 1 & 227 & 34 (15.0\%) & 1.163 \\
\bottomrule
\end{tabular}
\begin{tablenotes}\footnotesize
\item The maximum batch size reached is two on MK01 and three on MK07 and
      MK10. Under a first-in-first-out travel rule, which collects more
      aggressively, the rates rise to 3 of 64 pick-ups on MK01, 13 of 99 on
      MK07, and 55 of 204 on MK10.
\end{tablenotes}
\end{threeparttable}
\end{table}

MK01, the only instance on which policies are trained here, is the instance
with the fewest batching opportunities: 1.4\% of pick-ups against 12.3\% on
MK10. A mechanism that only acts when a batch can be formed will therefore
appear inert on MK01 for a reason that belongs to the instance. The
\texttt{m-no-batch} row of Table~\ref{tab:ablation} should be read with that in
mind, and it should not be read as a verdict on batching.

\subsection{Advantage ablation}
\label{sec:results:advantage}

This is the one experiment in the paper with five training seeds per arm,
because its conclusion is about the training procedure itself.
Table~\ref{tab:advantage} reports it.

\begin{table}[htbp]
\centering
\caption{Advantage ablation, five training seeds per arm, budget mechanism off,
all other switches identical. $p$ values are paired $t$-tests over the five
training seeds.}
\label{tab:advantage}
\setlength{\tabcolsep}{4pt}
\begin{threeparttable}
\begin{tabular}{@{}lrrr@{}}
\toprule
 & \texttt{scalar} & \texttt{reinforce} & \texttt{raw} \\
 & $(r-\bar r)/\sigma$ & $r - \bar r$ & $r$ \\
\midrule
Seed 0 & 84.04 & 71.12 & 84.97 \\
Seed 1 & 81.24 & 79.96 & 87.98 \\
Seed 2 & 87.19 & 75.31 & 90.72 \\
Seed 3 & 79.27 & 71.79 & 79.35 \\
Seed 4 & 71.56 & 85.85 & 85.41 \\
\midrule
Mean makespan & $80.66 \pm 5.90$ & $76.81 \pm 6.15$ & $85.69 \pm 4.22$ \\
Mean energy & 6.315 & 6.053 & 6.650 \\
Mean \TWT & 17.10 & 16.94 & 40.14 \\
\bottomrule
\end{tabular}
\begin{tablenotes}\footnotesize
\item Paired differences, makespan: \texttt{reinforce} $-$ \texttt{scalar}
      $= -3.85$ ($p = 0.48$); \texttt{raw} $-$ \texttt{scalar} $= +5.03$
      ($p = 0.11$); \texttt{raw} $-$ \texttt{reinforce} $= +8.88$
      ($p = 0.034$). Nine paired tests are reported in this section; a
      Bonferroni threshold over them is $0.0056$, which none of these reaches.
\end{tablenotes}
\end{threeparttable}
\end{table}

\begin{figure}[htbp]
\centering
\includegraphics[width=\columnwidth]{fig_advantage.pdf}
\caption{Advantage ablation, five training seeds per arm. Points are individual
seeds; the horizontal bar is the arm mean and the vertical bar is one standard
deviation. The three arms are identical in every switch and hyperparameter;
only the advantage differs.}
\label{fig:advantage}
\end{figure}

Figure~\ref{fig:advantage} gives two findings.

Subtracting a group-mean baseline has an effect. The arm that uses the raw
reward, with no baseline subtracted, is worse than the arm that subtracts the
group mean by 8.88 makespan units, with four of five seeds agreeing on the
sign. The uncorrected $p$ value is 0.034; with nine paired tests in this
section, the Bonferroni threshold is 0.0056, so we report this as a consistent
direction rather than as a significant effect.

Dividing by the group standard deviation has no measurable effect. The
\texttt{scalar} and \texttt{reinforce} arms differ by 3.85 makespan units with
$p = 0.48$, and the per-seed differences change sign. On this problem the
normalization step of the group-relative advantage is not doing detectable
work; the baseline subtraction is.

This distinction matters for how the method is described. The usual statement
is that group-relative advantages are necessary. What the data support is
narrower: a group-mean baseline is necessary, and the division by the group
standard deviation is not measurably different from omitting it.

This table is also where the seed variance was measured, and it is the reason
the rest of this paper's ablations are labelled as they are. The per-arm
standard deviation is 4--6 makespan units; the differences between arms are
4--9. Before this experiment was run, a single-seed comparison of the
\texttt{scalar} and \texttt{reinforce} arms read as a 15.4\% win for
\texttt{reinforce}; with five seeds it is 4.8\% with $p=0.48$. The
\texttt{scalar} arm's own seed 4 (71.56) is lower than the
\texttt{reinforce} arm's seed 0 (71.12) for all practical purposes.

\subsection{The objectives}
\label{sec:results:front}

Table~\ref{tab:front} reports five preference vectors trained to completion,
plus the default weight vector derived from the reference schedule.

\begin{table}[htbp]
\centering
\caption{Preference sweep on tier B. The three
one-hot vertices were trained as the single-objective baseline arms of
Table~\ref{tab:main}. ``Non-dominated'' is over the six rows of this table.}
\label{tab:front}
\setlength{\tabcolsep}{4pt}
\begin{threeparttable}
\begin{tabular}{@{}lrrrr@{}}
\toprule
Point & Preferences & Makespan & Energy & \TWT \\
\midrule
\texttt{b-ms} & $(1,0,0)$ & \textbf{74.09} & \textbf{5.790} & 5.977 \\
\texttt{b-twt} & $(0,0,1)$ & 77.30 & 6.048 & 7.392 \\
\texttt{e-me} & $(0.5,0.5,0)$ & 77.27 & 6.053 & 9.594 \\
\texttt{full} & $w$-derived & 78.90 & 6.160 & \textbf{5.042} \\
\texttt{e-eq} & $(1/3,1/3,1/3)$ & 80.96 & 6.333 & 20.456 \\
\texttt{b-en} & $(0,1,0)$ & 96.93 & 7.434 & 97.458 \\
\bottomrule
\end{tabular}
\begin{tablenotes}\footnotesize
\item The default weight vector is $w = (0.062, 0.819, 0.120)$, derived from the
      reference schedule. ``Non-dominated'' marks \texttt{b-ms} and
      \texttt{full}. All six runs completed all ten jobs with no horizon
      truncation.
\end{tablenotes}
\end{threeparttable}
\end{table}

The attainable front collapses to two points: \texttt{b-ms} and \texttt{full}.
The three remaining preference vectors, including both interior points, are
dominated.

The collapse is a property of the objectives, not of the search. On this
instance the three objectives are strongly aligned: a policy that does well on
makespan also does well on energy and tardiness. The makespan-only arm is not
merely non-dominated; it is within 0.9 tardiness units of the default arm while
being better on the other two objectives.

Two caveats. The endpoint of a single training run moves: the
\texttt{e-me} arm read 72.65 / 5.683 / 3.038 at step 270, which would have
dominated every point in the table, and 77.27 / 6.053 / 9.594 at step 300. Five
evaluation seeds do not pin the endpoint of a single arm. And the effective
strength of the budget mechanism differs between preference arms, because the
penalty enters before group normalization and is therefore scaled by the group
reward spread; the arms in Table~\ref{tab:front} differ in both the preference
vector and the effective penalty strength.

\subsubsection{The energy account and the makespan}
\label{sec:results:energy}

The paper reports energy as a third objective, and one property of the account
has to be stated so that the energy column is not read as something it is not.
Total energy as defined by Equation~\eqref{eq:energy} is close to
rank-equivalent with the makespan: across the thirteen ablation runs the ratio
of the two spans 7.6\%, and their rank correlation is $\rho = 0.978$.

The reason is not a finding but a consequence of the definition. Machine
processing accounts for 38--45\% of the total, machine standby for 21--31\%,
machine setup for 13--17\%, vehicles for 8--9\%, and the fixed shop load for
8.5\%; of these, the fixed shop load is charged over the makespan explicitly,
vehicle standby is likewise horizon-bound, and machine standby is not merely
correlated with the makespan but determined by it, since the time a machine
spends idle is the horizon minus the time it spends working. Together these
terms are a little under half of the total, and they are affine in the makespan
for any schedule whatsoever. The aggregate account therefore cannot move
independently of the makespan by more than the variation in the remaining
terms, and a schedule cannot be made more energy-efficient under this account
without also being made shorter. A net account that keeps only the terms
attributable to work actually performed, and drops the three horizon-bound
ones, is not tied to the makespan in this way ($\rho = 0.582$); it is reported
in Equation~\eqref{eq:energy:net} and used in Section~\ref{sec:results:front},
but the reward is trained on the total account, so it is a reading and not an
objective here.

Two consequences for reading the tables. The energy column of
Table~\ref{tab:main} should be read as a restatement of the makespan rather
than as an independent axis: under the standard account the reported
three-objective problem has an effective dimension of two. And the same is true
of the published work this paper compares against indirectly, because the
account is the standard one; the point is about the account, not about any
particular study. Section~\ref{sec:conclusion} records what would be needed to
turn the net account into an objective.

A counterfactual bounds the sensitivity to one modelling choice. The machine
standby power is taken from the loaded-idle column of the cited source rather
than its separately tabulated standby power; replacing it with the lower
tabulated value reduces the mean total energy by 12.9\% and widens the ratio
spread from 7.6\% to 16.6\%, while leaving the rank correlation unchanged at
$\rho = 0.973$. The choice moves the level of the reported energy, not its
ordering.

\subsubsection{Sensitivity to the assumed parameters}
\label{sec:results:sensitivity}

Most constraint parameters have no published source, and the sensitivity of the
reported results to them is part of what the paper has to establish.
Table~\ref{tab:sens} re-evaluates the tier B policy under perturbed dynamics.

\begin{table*}[htbp]
\centering
\caption{Sensitivity of the tier B policy to the assumed parameters, five
evaluation seeds. The policy is held fixed and only the dynamics change. The
default column repeats the tier B row of Table~\ref{tab:main}. Values are
makespan / energy / \TWT.}
\label{tab:sens}
\setlength{\tabcolsep}{5pt}
\begin{threeparttable}
\small
\begin{tabular}{@{}lrrr@{}}
\toprule
Parameter & Low & Default & High \\
\midrule
Setup time (min) & 1.0 & 2.0 & 4.0 \\
\quad & 67.38 / 5.359 / 0.00 & 78.90 / 6.160 / 5.04 & 95.73 / 7.315 / 100.07 \\
Failure rate & $\times 0.5$ & $\times 1$ & $\times 2$ \\
\quad & 78.81 / 6.152 / 4.89 & 78.90 / 6.160 / 5.04 & 84.74 / 6.593 / 45.41 \\
Rework probability & 0.025 & 0.05 & 0.10 \\
\quad & 73.53 / 5.773 / 3.11 & 78.90 / 6.160 / 5.04 & 81.99 / 6.399 / 8.05 \\
Maintenance interval (min) & 240 & 120 & 60 \\
\quad & \multicolumn{3}{c}{bit-identical to default at 240} \\
\quad & --- & 78.90 / 6.160 / 5.04 & 79.90 / 6.229 / 5.04 \\
Vehicle MTBF (min) & 240 & 480 & 960 \\
\quad & 80.02 / 6.240 / 7.73 & 78.90 / 6.160 / 5.04 & 78.81 / 6.154 / 5.16 \\
\bottomrule
\end{tabular}
\begin{tablenotes}\footnotesize
\item Job weights are not swept: the evaluation path hard-codes equal weights,
      so the parameter has no entry point. Re-evaluation holds the policy and
      its normalization context fixed and changes only the dynamics. The
      maintenance row is identical at 240 because MK01 never reaches the
      default interval in the first place.
\end{tablenotes}
\end{threeparttable}
\end{table*}

\begin{figure}[htbp]
\centering
\includegraphics[width=\columnwidth]{fig_sensitivity.pdf}
\caption{Sensitivity to the assumed parameters, the same numbers as
Table~\ref{tab:sens}. The policy is held fixed and only the dynamics change.
Parameter levels are not commensurate, so the horizontal axis is the level
index rather than a physical value, and each curve is identified by marker and
line style as well as by colour. The setup time and the failure rate dominate
the makespan; the tardiness panel separates the parameters far more sharply,
and the maintenance interval is flat on both because the instance never
reaches it.}
\label{fig:sensitivity}
\end{figure}

The two largest levers, in Figure~\ref{fig:sensitivity}, are the setup time and
the machine failure rate. Over
the swept range the setup time moves the makespan by 28 points and the
tardiness by 100, and the failure rate by 5.8 and 40. Rework is intermediate,
at 8.5 and 5. Vehicle failure and the maintenance interval move little, and the
maintenance parameter at 240 is bit-identical to the default because the
instance never reaches the interval at either setting.

Re-evaluation approximates retraining for these parameters. We trained two
fresh policies in the perturbed environment, at the most adverse level of each
of the two largest levers. For the setup time the retrained policy reaches
102.85 against 95.73 for the re-evaluated one; for the failure rate, 82.91
against 84.74. Both gaps are of the same order as the training-seed standard
deviation of 5.90 for makespan, so the order the re-evaluation reports is
informative. Tardiness moves much further, by $+18.1$ and $-24.5$, and the
re-evaluated tardiness should not be read as a retrained value.

\subsection{Discussion}
\label{sec:discussion}

The two tiers say more together than either says alone. On tier B, where the
ten constraints are present, the learned policy dominates every out-of-sample
point of the NSGA-II front and beats the rule baseline by 26.5 in makespan. On
tier A, where the constraints are switched off and the problem is easy, the
ordering reverses and the metaheuristic edges ahead. That pair locates the
policy's advantage in its representation of the constrained state rather than
in its search behaviour, which is the claim this paper set out to test.

The advantage ablation narrows a claim that is often stated too broadly. A
group-mean baseline is doing measurable work; the division by the group
standard deviation is not measurably different from omitting it. That is a
statement about the training procedure and not about this problem in
particular, and it is the one comparison in the paper that rests on enough
training seeds to be read as an estimate rather than as a measurement.

The constraint results are consistent rather than surprising. Every mechanism
that the instance can exercise shows a positive difference when it is removed,
and every mechanism that it cannot exercise sits at zero; the two groups are
separated by properties of the constraint that were measured independently of
the policy. What the ablation does not support is an ordering among the
mechanisms, and the reason is sample size rather than effect size.

